\chapter{Software Architecture}

\section*{Purpose}
This chapter defines the current Linux-side software structure hosted on the Zynq PS.

\section{Current Implemented Structure}
The repository contains both a deterministic local development path and a target-facing real-MMIO bring-up path. The public software contract remains unchanged: the daemon accepts a 64-byte digest, sequences the MMIO device abstraction, and returns signature bytes plus status.

\section{Current Components}
\begin{itemize}[leftmargin=*]
  \item \textbf{Service core:} transport-independent signing orchestration over the device abstraction.
  \item \textbf{MMIO device abstraction:} stable software API for digest writes, start, wait, readback, and clear-status behavior.
  \item \textbf{Fake backend:} deterministic local register model used for unit tests, local self-test, and client smoke checks.
  \item \textbf{Real backend:} PoC user-space MMIO backend intended for Linux on Zynq, using a `/dev/mem`-style mapping path and the documented wrapper base address.
  \item \textbf{Probe helper:} safe register-visibility path that does not trigger signing.
  \item \textbf{STUB selftest path:} explicit one-transaction bring-up path that validates the deterministic STUB signature rule.
\end{itemize}

\section{Backend Selection}
Backend selection is configuration-driven. The default remains `fake`, so local development continues to work exactly as before. The `real` backend is selected through environment variables or command-line overrides and requires platform-specific MMIO mapping inputs such as the wrapper base address.

\section{Current Bring-Up Role}
The software stack is now prepared for the first real board interaction in STUB mode. The real probe command checks register visibility, and the real selftest writes a known digest, starts one operation, waits for completion, reads the signature, and compares it against the documented deterministic STUB rule.

\section{Current Claim Boundary}
The repository now contains the software infrastructure required for first PS-to-PL bring-up in STUB mode, but it does not yet claim that target-board register access or STUB selftest has been validated on actual hardware.